\makeabstract
    {
    Optimization of Live-dead Characterization of Mycobacterium avium subsp. hominissuis subpopulation
    }
    {
    Adrita Zaima Islam\textsuperscript{1}, India Ingemi\textsuperscript{2}, Sloan Siegrist\textsuperscript{2}
    }
    {
    \textsuperscript{1} Amherst College, Amherst, MA\\
\textsuperscript{2} Department of Microbiology, University of Massachusetts Amherst, Amherst, MA
    }
    {
    The CHAMP assay provides a method for investigating the accumulation of molecules within the bacterial cytosol in bacteria that are modified to express the HaloTag (HT) protein. Adapting the assay for use in mycobacterial species, known for an outer membrane rich in mycolic acids that hinders the permeation of molecules, including drugs, will be crucial for furthering drug discovery efforts. During initial adaptation of the assay in Mycobacterium avium, the most common cause of non-tuberculous mycobacterial (NTM) infections in humans, the fluorophore CF-647-Az, previously selected for use in the PAC-MAN assay, which queries permeability of molecules through the peptidoglycan layer of mycobacteria, was unexpectedly restricted to staining approximately 3\% of the bacterial population. We were curious whether this 3\% was representative of the entire MAH11 population, or whether this subpopulation had a unique permeability phenotype, either due to being dead cells, or due to differential cell envelope permeability, either driven phenotypically or genetically.
    }
    {
    To address this question, we developed a multiparameter flow cytometric live/dead triple staining approach for M. avium subsp. hominissuis strain 11 (MAH11). Heat-killed cells were established as a dead-cell control. Flow cytometry was used to optimize the concentration and incubation time of propidium iodide (PI), a membrane-impermeant viability dye. Similarly, 7-hydroxycoumarin-carbonylamino-D-alanine (HADA) was evaluated as a viability-associated stain. These conditions were incorporated into a triple-staining workflow combining PI, CF-647-Az, and HADA. Finally, the results of the staining protocol were evaluated using flow cytometry.  
    }
    {
    HADA stained only {\textasciitilde}52{\textendash}55\% of non-heat killed cells while also labeling {\textasciitilde}8{\textendash}16\% of the heat-killed cells, limiting its ability to reliably distinguish live from dead M. avium. Using a different live stain or a different HADA sample will be part of a future study. PI similarly stained {\textasciitilde}55\% of the non-heat-killed cells, suggesting either excessive membrane penetration due to too high a concentration of PI or substantial cell death. Although CF-647-Az labeled the expected {\textasciitilde}3\% population in non-heat-killed cells, it also labeled {\textasciitilde}60{\textendash}66\% of the heat-killed cells, suggesting that membrane compromise may contribute to labeling. Additionally, MAH11 which was being grown in 7H9+ADC, showed slowed growth, suggesting compromised health of the bacterial cells. Growing MAH11 in 7H9+OADC saw a significant improvement in bacterial growth, suggesting that the MAH11 population we were querying was not receiving adequate nutrients for typical growth and may suggest the possibility of dead or dying cells with a compromised cell envelope. 
    }
    {
    Overall, the triple-stain results were inconclusive. Repetition of the assay under the improved growth conditions in 7H9+OADC may be helpful.
    }

    \newpage
    